\section{Conclusion}
\label{sec:conclusion}

In this paper, we investigated the "Bullying Paradox": the superior effectiveness of iterative rejection over single-shot instructions for AI detection evasion. We demonstrated that five rounds of "bullying" an LLM can nearly double the human-ness score of its output compared to a direct "be human" prompt. This improvement is linked to a qualitative regime shift in the model's linguistic style, characterized by higher perplexity and a rebound in sentence length variation.

Our findings suggest that negative constraints provide a more powerful steering signal for escaping the AI latent plateau than positive goals. As AI-generated text becomes more pervasive, understanding these conversational vulnerabilities is critical for developing more robust detection and safety systems. Future work will focus on scaling this analysis to larger datasets and probing the internal activations of models under iterative pressure.
